% !TeX root = 00-main.tex
\documentclass[00-main.tex]{subfiles}

\begin{document}

\section{Introduction}
The classical \textit{susceptible--infectious--recovered} (SIR) model provides one of the simplest mathematical frameworks for describing the spread of an infectious disease through a population. At time $t$, let $\widehat S(t)$, $\widehat I(t)$, and $\widehat R(t)$ denote the numbers of susceptible, infectious, and recovered or removed individuals, respectively. The total population is therefore
\[
    N(t)=\widehat S(t)+\widehat I(t)+\widehat R(t).
\]
We normalize the compartment sizes by $N(t)$ and, for simplicity, reuse the letters $S$, $I$, and $R$ for the resulting population proportions:
\[
    S(t)=\frac{\widehat S(t)}{N(t)}, \qquad
    I(t)=\frac{\widehat I(t)}{N(t)}, \qquad
    R(t)=\frac{\widehat R(t)}{N(t)}.
\]
Consequently, $S(t)+I(t)+R(t)=1$, and each state variable takes values in $[0,1]$.
To allow for demographic turnover, we consider the normalized SIR system
\[
\begin{cases}
		S'(t) = \Lambda - \beta S(t)I(t) - \mu S(t), \\
		I'(t) = \beta S(t)I(t) - \gamma I(t) - \mu I(t), \\
		R'(t) = \gamma I(t) - \mu R(t).
\end{cases}
\]
Here $\Lambda$ is the \textit{recruitment rate} into the susceptible compartment per unit population, $\beta$ is the \textit{effective transmission rate}, $\gamma$ is the \textit{recovery/removal rate}, and $\mu$ is the \textit{natural per-capita mortality rate}.

The nonlinear term $\beta SI$ describes the incidence of new infections under homogeneous mixing: susceptible individuals leave the $S$ compartment and enter the $I$ compartment at this rate.
Infectious individuals recover or are removed at rate $\gamma I$, while natural mortality acts on all three compartments at rate $\mu$.

Adding the three equations gives
\[
    \frac{d}{dt}(S+I+R)=\Lambda-\mu(S+I+R).
\]
Thus, under the normalization $S+I+R=1$, demographic equilibrium requires $\Lambda=\mu$; this assumption keeps the normalized population constant.
The special case $\Lambda=\mu=0$ reduces the system to the classical closed-population SIR model.

Six principal references determine the scope and organization of this note.
The review in \cite{Ch21} provides the broad collection of fractional compartmental models summarized in Section~2, while \cite{Sa17} supplies the memory-kernel construction used in this section to motivate the transition from a classical SIR system to a fractional one.
The deterministic and stochastic SEIR comparison in Section~3 is based on \cite{Is20}, which offers a focused case study of what changes when an integer-order epidemic model is replaced by a fractional-order model.
The final critical discussion draws on three complementary perspectives: \cite{Gh25} examines dimensional consistency and parameter interpretation, \cite{Cu26} questions the biological meaning assigned to fractional memory, and \cite{Ci26} emphasizes theoretical construction, estimation, and parameter-identification difficulties.
Together, these sources provide the introductory, comparative, and critical viewpoints selected for this reading note; they are not intended to constitute an exhaustive survey of the literature.


\subsection{Limitations of the classical SIR model}

Although the classical SIR model provides a useful first approximation of epidemic dynamics, it relies on several simplifying assumptions:
\begin{enumerate}
	\item \textit{Lack of explicit memory.}
	The system is memoryless in the sense that its future evolution is determined entirely by the current values of $S(t)$, $I(t)$, and $R(t)$.
	The rates of infection and recovery at time $t$ do not explicitly depend on the previous states of the epidemic.
	Consequently, the model may not adequately represent processes with memory or delay, such as variable incubation periods, heterogeneous recovery times, temporary immunity, or behavioral responses influenced by past infection levels.

	\item \textit{Constant parameters.}
	The parameters $\lambda$, $\beta$, $\gamma$, and $\mu$ are usually assumed to be constant over time.
	In practice, however, the transmission rate may change because of public-health interventions, seasonal effects, changes in contact behavior, or the emergence of new variants.
	Similarly, recovery and mortality rates may vary with medical treatment, healthcare capacity, demographic characteristics, and the stage of the epidemic.
	Treating these parameters as constants therefore limits the ability of the classical model to capture time-dependent and history-dependent epidemic behavior.
\end{enumerate}

These limitations motivate extensions of the classical SIR framework.
One possible approach is to replace the ordinary time derivatives with \textit{fractional derivatives}, which introduce nonlocal dependence on the past states of the system and provide a mathematical mechanism for modeling memory effects.


\subsection{Preliminaries}

In this section, we introduce the basic concepts of fractional calculus needed for the fractional SIR models discussed later.
We first motivate the fractional integral through the Cauchy formula for repeated integration, and then present the Caputo and Riemann--Liouville fractional derivatives together with the relation between them.

The definition of a fractional integral is motivated by the classical formula for repeated integration. Let
\[
    {}_a I_t f(t) = \int_a^t f(\tau)\,d\tau
\]
denote the integral of $f$ from $a$ to $t$. Applying this operator $n$ times gives the $n$-fold integral
\[
    {}_a I_t^n f(t)
    = \int_a^t \int_a^{t_1} \cdots \int_a^{t_{n-1}}
      f(t_n)\,dt_n\cdots dt_2\,dt_1.
\]
By changing the order of integration, this iterated integral can be written as a single integral. The resulting identity is the \emph{Cauchy formula for repeated integration}:
\[
    {}_a I_t^n f(t)
    = \frac{1}{(n-1)!}\int_a^t (t-\tau)^{n-1}f(\tau)\,d\tau,
    \qquad n\in\mathbb{N}.
\]
The right-hand side depends on $n$ through the power $n-1$ and the factorial $(n-1)!$. Since the gamma function satisfies
\[
    \Gamma(z)=\int_0^\infty s^{z-1}e^{-s}\,ds,
    \qquad \Gamma(n)=(n-1)!,
\]
the integer order $n$ can be replaced naturally by an arbitrary real order $\alpha>0$. This observation leads to the following definition.

\begin{definition}[Fractional integral]
	The \textit{fractional integral} of order $\alpha>0$ of a function $f$ is defined by
	\[
        {}_a I_t^\alpha f(t)
        = {}_a D_t^{-\alpha}f(t)
        = \frac{1}{\Gamma(\alpha)}
          \int_a^t (t-\tau)^{\alpha-1}f(\tau)\,d\tau.
    \]
\end{definition}

Unlike fractional integration, fractional differentiation does not have a unique standard definition.
Two commonly used definitions are the \textit{Caputo derivatives} and \textit{Riemann--Liouville derivatives}.
The Caputo derivative first applies an integer-order derivative and then a fractional integral, whereas the Riemann--Liouville derivative reverses this order by applying a fractional integral before the integer-order derivative.

\begin{definition}[Caputo derivative]
	The \textit{Caputo derivative} of order $\alpha$ of a function $f(t)$ is defined by
	\[ {}_{a}^{\mathrm{C}}\!D_t^\alpha f(t)
	= {}_{a}\!D_t^{-(n-\alpha)} \frac{d^n}{dt^n} f(t)
	= \frac{1}{\Gamma(n-\alpha)} \int_a^t (t-\tau)^{n-\alpha-1} f^{(n)}(\tau)\,d\tau,
	\]
	where $n-1 < \alpha < n$ and $n\in\mathbb{Z}^+$.
\end{definition}

\begin{definition}[Riemann--Liouville derivative]
	The \textit{Riemann--Liouville derivative} of order $\alpha$ of a function $f(t)$ is defined by
	\[ {}_{\phantom{R}a}^{\mathrm{RL}}\!D_t^\alpha f(t)
	= \frac{d^n}{dt^n} {}_{a}\!D_t^{-(n-\alpha)} f(t)
	= \frac{1}{\Gamma(n-\alpha)} \frac{d^n}{dt^n} \int_a^t (t-\tau)^{n-\alpha-1} f(\tau)\,d\tau,
	\]
	where $n-1 < \alpha < n$ and $n\in\mathbb{Z}^+$.
\end{definition}

These two derivatives are not equivalent in general.
In fact, their relation can be described as:
\[ {}_a^\mathrm{C}\!D_t^\alpha f(t) = {}_{\phantom{R}a}^\mathrm{RL}\!D_t^\alpha f(t) - \sum\limits_{k=0}^{n-1} \frac{(t-a)^{k-\alpha} f^{(k)}(a)}{\Gamma(k-\alpha+1)}. \]
If the initial conditions $f^{(k)}(a) = 0$ for all $k\in\{0,1,\ldots,n-1\}$, then they are equivalent.


\subsection{Fractional SIR model}

Having identified the absence of explicit memory as a limitation of the classical model, we now make the desired history dependence mathematically precise.
We follow the kernel-based construction of \cite{Sa17}, beginning for clarity with the closed-population SIR model:
\[
\begin{cases}
    S'(t)=-\beta S(t)I(t), \\
    I'(t)=\beta S(t)I(t)-\gamma I(t), \\
    R'(t)=\gamma I(t).
\end{cases}
\]
Following \cite{Sa17}, memory is introduced by allowing the present rates of change to depend on transition rates at all earlier times.
This gives the convolution system
\[
\begin{aligned}
    \frac{dS(t)}{dt}
        &=-\beta\int_0^t
          \kappa(t-\tau)S(\tau)I(\tau)\,d\tau, \\
    \frac{dI(t)}{dt}
        &=\int_0^t\kappa(t-\tau)
          \bigl[\beta S(\tau)I(\tau)
          -\gamma I(\tau)\bigr]\,d\tau, \\
    \frac{dR(t)}{dt}
        &=\gamma\int_0^t
          \kappa(t-\tau)I(\tau)\,d\tau,
\end{aligned}
\]
where $\kappa(t-\tau)$ is a memory kernel.
In the Markovian limit, the kernel is concentrated at the present time and is represented by a Dirac delta function.
To model long-range memory, \cite{Sa17} choose the power-law kernel
\[
    \kappa(t-\tau)
    =\frac{(t-\tau)^{\alpha-2}}{\Gamma(\alpha-1)}.
\]
The occurrence of $\Gamma(\alpha-1)$, rather than $\Gamma(\alpha)$, results from differentiating the corresponding fractional integral:
\[
    \frac{d}{dt}
    \left[
        \frac{1}{\Gamma(\alpha)}
        \int_0^t(t-\tau)^{\alpha-1}F(\tau)\,d\tau
    \right]
    =
    \frac{1}{\Gamma(\alpha-1)}
    \int_0^t(t-\tau)^{\alpha-2}F(\tau)\,d\tau.
\]
This identity is understood in the fractional-operator sense; for $0<\alpha<1$, the expression on the right is hypersingular and is not an ordinary improper integral. Rewriting the convolution system in fractional-operator notation gives the Caputo SIR model
\[
\begin{cases}
    {}_{0}^{\mathrm{C}}\!D_t^\alpha S(t)
        =-\beta^\alpha S(t)I(t), \\
    {}_{0}^{\mathrm{C}}\!D_t^\alpha I(t)
        =\beta^\alpha S(t)I(t)-\gamma^\alpha I(t), \\
    {}_{0}^{\mathrm{C}}\!D_t^\alpha R(t)
        =\gamma^\alpha I(t),
\end{cases}
\]
subject to the initial conditions
\[
    S(0)=S_0, \qquad I(0)=I_0, \qquad R(0)=R_0,
    \qquad S_0+I_0+R_0=1.
\]
The superscript $\alpha$ distinguishes the fractional-order parameters $\beta^\alpha$ and $\gamma^\alpha$ from their classical counterparts.
For dimensional consistency, the coefficients appearing in this fractional system must have units compatible with $\mathrm{time}^{-\alpha}$.
However, interpreting them simply as powers of the classical rates is itself a modeling choice; alternative dimensional corrections have also been proposed in the literature and will be discussed later.

As $\alpha\to 1$, the memory kernel approaches the memoryless limit and the Caputo derivative reduces to the ordinary first derivative, so the classical SIR model is recovered.
For $0<\alpha<1$, the power-law memory has no finite cutoff: every earlier state contributes to the present dynamics, although with a time-dependent weight.
\cite{Sa17} shows that this memory mechanism can alter epidemic thresholds and transient behavior in both fully mixed populations and structured networks.


\subsection{Why use a fractional model?}

The preceding construction gives a precise mathematical route from a memory kernel to a fractional SIR model, but it does not imply that fractional differentiation is the only possible description of epidemic memory.
The choice should be evaluated by considering both its advantages and its limitations.
\begin{enumerate}
    \item \textit{Arguments in favor.}
    A Caputo derivative is nonlocal, so the current rate of change depends explicitly on the past trajectory rather than only on the present compartment sizes.
    Its power-law kernel provides a parsimonious representation of slowly decaying memory using a single additional parameter $\alpha$, while retaining the classical model as the limiting case $\alpha=1$.
    This makes it possible to study how memory changes transient dynamics and epidemic thresholds within a framework closely related to the ordinary SIR model.

    \item \textit{Points requiring caution.}
    Replacing an ordinary derivative by a fractional one is a modeling assumption, not a consequence of non-Markovian behavior alone.
    The power-law kernel should be supported by a plausible mechanism or by data; otherwise, $\alpha$ may serve only as an extra fitting parameter without a clear epidemiological interpretation.
    Fractional parameters also require careful dimensional treatment and may be difficult to identify separately from transmission, recovery, and initial-condition effects.

    \item \textit{Alternative descriptions.}
    Memory can also be modeled through delay differential equations, renewal equations, infection-age or stage-structured models, distributed waiting times, and explicitly time-dependent parameters.
    These alternatives may be preferable when the relevant mechanism is a specific incubation delay, a known infectious-period distribution, or a documented intervention schedule rather than scale-free long-range memory.
\end{enumerate}

Therefore, the main justification for using a fractional SIR model is not merely that real epidemics have a past.
It is that the epidemic exhibits history dependence that can reasonably be represented by a nonlocal power-law kernel, and that the fractional model provides a better balance of explanatory value, predictive performance, and parsimony than suitable non-fractional alternatives.
In applications, this claim should ultimately be assessed by comparing competing models against data rather than assumed in advance.

\ifSubfilesClassLoaded{\printbibliography}{}

\end{document}
